For multifold sums of powers, we use the notation provided by Donald Knuth in~\cite{knuth1993johann},
because it is simple and beautiful.
\begin{definition} [Multifold sums of powers]
  \label{def:multifold-sums-of-powers-recurrence}
  For non-negative integers $r,n,m$
  \begin{align*}
    \KnuthRFoldSum{0}{n}{m}   &= n^m, \\
    \KnuthRFoldSum{1}{n}{m}   &= \KnuthRFoldSum{0}{1}{m}
    + \KnuthRFoldSum{0}{2}{m} + \cdots + \KnuthRFoldSum{0}{n}{m}, \\
    \KnuthRFoldSum{r+1}{n}{m} &= \KnuthRFoldSum{r}{1}{m}
    + \KnuthRFoldSum{r}{2}{m} + \cdots + \KnuthRFoldSum{r}{n}{m}.
  \end{align*}
\end{definition}

We utilize the following notation for rising factorials
\begin{definition}[Rising factorials]
  \label{def:rising-factorials}
  For integers $n,k$
  \begin{align*}
    \risingFactorial{n}{k} =
    \begin{cases}
      0, \quad & \mathrm{if \; } k < 0, \\
      1, \quad & \mathrm{if \; } k = 0, \\
      \prod_{j=0}^{k-1} (n+j), \quad & \mathrm{if \; } k>0.
    \end{cases}
  \end{align*}
\end{definition}
We encourage the audience to read
J. F. Steffensen's work on the definition of
generalized factorials~\cite{steffensen1933definition}.

\begin{lemma}[Forward Newton's formula]
  \label{lem:forward-newtons-formula}
  Let $f\colon \mathbb{Z} \to \mathbb{C}$ be a function,
  and let $t \in \mathbb{Z}$.
  Then, for all $n \in \mathbb{Z}$,
  \begin{align*}
    f(n)
    = \sum_{k=0}^{\infty} \frac{\risingFactorial{(n-t)}{k}}{k!} \Delta^{k} f(t)
    = \sum_{k=0}^{\infty} \binom{n-t}{k} \Delta^{k} f(t),
  \end{align*}
  where the $k$-th forward difference $\Delta^{k} f(t)$ of $f$
  evaluated at $t$ is defined by
  \begin{align*}
    \Delta^{k} f(t)
    = \tsum_{j=0}^{k} (-1)^{k-j} \tbinom{k}{j} f(t+j).
  \end{align*}
  \begin{proof}
    By~\cite[section 2, eq. (19)]{steffensen1927interpolation},
    and by identity $\frac{\risingFactorial{(n-t)}{k}}{k!} = \binom{n-t}{k}$.
  \end{proof}
\end{lemma}

\begin{lemma}[Forward Newton's formula for powers]
  \label{lem:forward-newtons-formula-for-powers}
  For non-negative integers $n,m$, and an arbitrary integer $t$
  \begin{align*}
    n^m = \sum_{k=0}^{m} \binom{n-t}{k} \Delta^{k} t^m.
  \end{align*}
  \begin{proof}
    Special case of Lemma~\eqref{lem:forward-newtons-formula} for powers.
  \end{proof}
\end{lemma}

\begin{lemma}[Generalized hockey-stick identity]
  \label{lem:generalized-hockey-stick-identity}
  For arbitrary integers $a,b$ and $j$,
  \begin{align*}
    \sum_{k=a}^{b} \binom{k}{j} = \binom{b+1}{j+1} - \binom{a}{j+1}.
  \end{align*}
  \begin{proof}
    We have,
    $\sum_{k=a}^{b} \binom{k}{j}
    = \binom{a}{j} + \binom{a+1}{j} + \cdots + \binom{b}{j}
    = \left( \sum_{k=0}^{b} \binom{k}{j} \right) - \left( \sum_{k=0}^{a-1} \binom{k}{j} \right)$.
    By hockey-stick identity $\sum_{k=0}^{n} \binom{k}{j} = \binom{n+1}{j+1}$ yields,
    \begin{align*}
      \tsum_{k=a}^{b} \tbinom{k}{j}
      = \left( \tsum_{k=0}^{b} \tbinom{k}{j} \right) - \left( \tsum_{k=0}^{a-1} \tbinom{k}{j} \right)
      = \tbinom{b+1}{j+1} - \tbinom{a}{j+1}.
    \end{align*}
    This completes the proof.
  \end{proof}
\end{lemma}

\begin{lemma}[Forward hockey-stick identity]
  \label{lem:forward-hockey-stick-identity}
  For integer $n\geq1$, and arbitrary integers $t,k,r$
  \begin{align*}
    \sum_{j=1}^{n} \binom{j-t+r}{k+r}
    = \binom{n-t+r+1}{k+r+1} - \binom{1-t+r}{k+r+1}.
  \end{align*}
  \begin{proof}
    By setting $a=1-t+r$, and $b=n-t+1$
    into Lemma~\eqref{lem:generalized-hockey-stick-identity}.
  \end{proof}
\end{lemma}

\begin{lemma}[Multifold sums of zero powers]
  \label{lem:multifold-sum-of-zero-powers}
  For non-negative integers $r,n$
  \begin{align*}
    \KnuthRFoldSum{r}{n}{0} = \binom{r+n-1}{r}.
  \end{align*}
  \begin{proof}
    \leavevmode\par
    \begin{enumerate}
      \item Let $r=0$, then we have $\KnuthRFoldSum{0}{n}{0} = 1$,
        by definition~\eqref{def:multifold-sums-of-powers-recurrence}.
      \item Let $r=1$, then we have
        $\KnuthRFoldSum{1}{n}{0}
        = \sum_{k=1}^{n} 1 = \binom{n}{1}$.
      \item Let $r=2$, then we have
        $\KnuthRFoldSum{2}{n}{0}
        = \sum_{k=1}^{n} \binom{k}{1}
        = \binom{n+1}{2}$.
      \item Let $r=3$, then we have
        $\KnuthRFoldSum{3}{n}{0}
        = \sum_{k=1}^{n} \binom{k+1}{2}
        = \binom{n+2}{3}$.
      \item Similarly, for arbitrary integer $r$, by hockey-stick identity
        $\sum_{k=1}^{n} \binom{k}{r} = \binom{n+1}{r+1}$,
        yields
        $\KnuthRFoldSum{r}{n}{0}
        =  \sum_{k=1}^{n} \KnuthRFoldSum{r-1}{k}{0}
        =  \sum_{k=1}^{n} \binom{k+r-2}{r-1}
        = \binom{r+n-1}{r}$.
    \end{enumerate}
    This completes the proof.
  \end{proof}
\end{lemma}

\begin{lemma}[Upper Binomial negation]
  \label{lem:upper-binomial-negation}
  For integers $n,k$
  \begin{align*}
    \binom{-n}{k} = (-1)^{k} \binom{k+n-1}{k}
  \end{align*}
  \begin{proof}
    See~\cite[p. 164, eq. (5.14)]{graham1994concrete}.
  \end{proof}
\end{lemma}

\begin{convention} For all $x$
  \begin{align*}
    x^0 = 1.
  \end{align*}
\end{convention}
Donald Knuth give extensive discussion on the convention $x^0=1$ for
all $x$, including zero, in Concrete Mathematics~\cite[~p. 162]{graham1994concrete},
and in~\cite{knuth1992notesnotation}.
